% Generated draft: requires mathematical and source review.
\documentclass[11pt]{article}
\usepackage[margin=25mm]{geometry}
\usepackage{fontspec}
\setmainfont{Latin Modern Roman}
\usepackage{amsmath,amssymb,mathtools,unicode-math}
\setmathfont{Latin Modern Math}
\usepackage{enumitem,needspace,xurl,hyperref}
\setlength{\parindent}{0pt}
\setlength{\parskip}{5pt}
\setlength{\emergencystretch}{3em}
\newcommand{\N}{\mathbb{N}}
\newcommand{\Z}{\mathbb{Z}}
\newcommand{\Q}{\mathbb{Q}}
\newcommand{\R}{\mathbb{R}}
\newcommand{\C}{\mathbb{C}}
\newcommand{\F}{\mathbb{F}}
\newcommand{\A}{\mathbb{A}}
\newcommand{\PP}{\mathbb{P}}
\newcommand{\E}{\mathbb{E}}
\newcommand{\eps}{\varepsilon}
\newcommand{\dd}{\,\mathrm{d}}
\begin{document}
\section*{UnsolvedMath ID: 10400014}
\textbf{Source problem code:} catalogue entry
\subsection{Definitions and mathematical statement}
For the knot invariant denoted $Q_K$ in the source, consider its restriction to the unit circle. In the intended polynomial/Laurent-polynomial setting define the maximizing set
\[
A(K)=\left\{\alpha\in\C:|\alpha|=1,\quad
|Q_K(\alpha)|=\max_{|z|=1}|Q_K(z)|\right\}.
\]
The source, attributed to A. Stoimenow, asks why these maximizing unit complex numbers are statistically concentrated around
\[
z_0=\exp(11\pi i/25).
\]
The collapsed printed expression is interpreted as this exponent; no different root or denominator is substituted.

The record does not define $Q_K$, its variable/normalization, or the knot population and sampling procedure. It must not be silently identified with the Jones polynomial merely because the surrounding subject is quantum invariants. The maximizing set can have multiple points; for a real-coefficient invariant complex conjugation gives symmetric modulus values, so choosing upper-half-circle representatives, counting all maximizers, or choosing one maximizer gives different empirical distributions.

Statistical concentration also lacks a size parameter, neighborhood/rate, or assertion of convergence. It could describe a finite experimental bias rather than asymptotic concentration at a point. Constant invariants and degenerate maxima need a convention. Thus the precise task is to explain or formulate the source's observed maximizing-angle phenomenon after those data are supplied, not an established universal assertion that every knot has its maximum at $z_0$.
\subsection{Short English statement}
Explain the reported concentration of unit-circle maximizers of $|Q_K|$ near $e^{11\pi i/25}$. The invariant, knot ensemble and rule for multiple/conjugate maxima are not specified.
\subsection{Sources}
\begin{itemize}
\item Local UnsolvedMath registry, permanent ID 10400014, exact source code catalogue entry, Problem 1.14; source attribution A. Stoimenow, with no proposer/year in metadata. Registry status open, classification OPEN-TRIAGE. The registry title is truncated; the mathematical question is retained from its statement.
\item Ohtsuki, \emph{Problems on invariants of knots and 3-manifolds} (2002), Problem 1.14, PDF page 19, \url{https://msp.org/gtm/2002/04/gtm-2002-04-024s.pdf}. Supplied pdf-font-structure extraction accessed July 29, 2026; provenance distinguishes 231 bold numbered items from surrounding references. No PDF was accessed here.
\item The supplied summary only describes a statistical concentration question for $Q_K$ at roots of unity and says it is not obviously resolved. It supplies no invariant definition, data set, theorem or later citation.
\item Provenance flags current primary-literature status and redistribution terms NEEDS\_REVIEW.
\end{itemize}
\end{document}
